\documentclass[10pt]{article}
\usepackage[T1]{fontenc}
\usepackage[utf8]{inputenc}
\usepackage[french]{babel}
\usepackage{amsmath,amssymb,mathrsfs,mathpazo}
\usepackage{geometry}
\geometry{a4paper,margin=2cm}
\usepackage{graphicx,array,tabularx,tabularray,colortbl,multirow,multicol,booktabs,enumitem,pifont,ulem,fancybox,diagbox,cancel}
\usepackage{pstricks,pstricks-add,pst-plot,pst-tree,pst-node,pst-text,pst-func,pst-math,pst-bspline,pst-3dplot,pst-eucl,pst-solides3d}
\let\pscalculate\undefined
\usepackage{pst-bezier}
\usepackage{tikz,pgfplots,tkz-tab}
\usetikzlibrary{arrows,arrows.meta,patterns,calc,positioning,trees,decorations.markings,shapes.geometric,shapes.misc,shapes.arrows,intersections,decorations.pathmorphing,decorations.pathreplacing,backgrounds,shadows,babel}
\pgfplotsset{compat=1.18}
\usepgfplotslibrary{fillbetween}
\usepackage[np]{numprint}
\usepackage{listings,makecell,pas-tableur,twoopt}
\let\barre\undefined
\usepackage[pstricks]{tablvar}
\definecolor{darkgreen}{RGB}{0,100,0}
\providecommand{\red}{\color{red}}
\providecommand{\blue}{\color{blue}}
\providecommand{\green}{\color{green}}
\providecommand{\black}{\color{black}}
\providecommand{\white}{\color{white}}
\providecommand{\gray}{\color{gray}}
\providecommand{\cyan}{\color{cyan}}
\providecommand{\magenta}{\color{magenta}}
\providecommand{\yellow}{\color{yellow}}
\providecommand{\violet}{\color{violet}}
\providecommand{\orange}{\color{orange}}
\providecommand{\euro}{\texteuro}
\providecommand{\EUR}{\texteuro}
\providecommand{\ieme}{\textsuperscript{e}}
\providecommand{\degres}{\ensuremath{^\circ}}
\providecommand{\degree}{\ensuremath{^\circ}}
\providecommand{\vv}[1]{\overrightarrow{#1}}
\providecommand{\Coord}[1]{#1}
\providecommand{\decosix}{}
\providecommand{\starredbullet}{\ensuremath{\bullet}}
\pagestyle{empty}
\setlength{\parindent}{0pt}
\usepackage{adjustbox}\DeclareUnicodeCharacter{2236}{:}\DeclareUnicodeCharacter{2190}{\ensuremath{\leftarrow}}
\providecommand{\R}{\mathbb{R}}
\providecommand{\N}{\mathbb{N}}
\providecommand{\Z}{\mathbb{Z}}
\providecommand{\Q}{\mathbb{Q}}
\providecommand{\D}{\mathbb{D}}
\providecommand{\C}{\mathbb{C}}
\providecommand{\e}{{\mathrm{e}}}
\providecommand{\vect}[1]{\overrightarrow{#1}}
\providecommand{\vectt}[1]{\overrightarrow{\mathrm{#1}}}
\providecommand{\vvt}[1]{\overrightarrow{\mathrm{#1}}}
\providecommand{\np}[2][]{{#2~#1}}
\providecommand{\vv}[1]{{\overrightarrow{#1}}}
\renewcommand{\d}{{\mathrm{d}}}
\providecommand{\degres}{{^\circ}}
\providecommand{\degre}{{^\circ}}
\providecommand{\blue}{\color{blue}}
\providecommand{\red}{\color{red}}
\providecommand{\enskip}{\quad}
\providecommand{\llbracket}{\lbrack\!\lbrack}
\providecommand{\rrbracket}{\rbrack\!\rbrack}
\providecommand{\euro}{{\text{€}}}
\providecommand{\no}{{\text{n°}}}
\providecommand{\interval}[2]{{[#1;#2]}}
\providecommand{\intervFO}[2]{{[#1;#2[}}
\providecommand{\dotfill}{\ldots}

\providecommand{\ds}{\displaystyle}
\providecommand{\textding}[1]{#1}
\providecommand{\Oij}{\(\mathrm{O};\vec{\imath},\vec{\jmath}\)}
\providecommand{\Oijk}{\(\mathrm{O};\vec{\imath},\vec{\jmath},\vec{k}\)}
\providecommand{\Ouv}{\(\mathrm{O};\vec{u},\vec{v}\)}
\providecommand{\intervFF}[2]{\ensuremath{[#1;#2]}}
\renewcommand{\euro}{\ifmmode\text{\texteuro}\else\texteuro\fi}
\DeclareUnicodeCharacter{00B0}{\ifmmode^{\circ}\else\textdegree\fi}
\DeclareUnicodeCharacter{20AC}{\ifmmode\text{\texteuro}\else\texteuro\fi}
% Énoncés : Éducation nationale. Transcription et corrigés : APMEP, Denis Vergès et contributeurs. Découpage et classement : CoopMaths, Sébastien Lozano, Éric Elter et contributeurs. Adaptation de la présentation : Maths Explorer.
% CC BY-SA : https://coopmaths.fr/www/about/
% https://coopmaths.fr/alea/static/bac/2022/tex/bac_2022_05_sujet2_madagascar_3_cor.tex
\begin{document}
\noindent\textit{APMEP / CoopMaths. CC BY-SA. Présentation adaptée par Maths Explorer.}\par\medskip
\medskip

\begin{enumerate}
	\item La fonction $g$ est continue et dérivable. $\forall x \in \R$, $f'(x)=\np{1000}x^{999}+1$.
	
La fonction $g'$ est continue et dérivable. $\forall x \in \R$, $g''(x)=999\times \np{1000}x^{998}=\np{999000}x^{998}$.
	
	$\forall x \in \R$, $\np{999000}x^{998}\geqslant 0$ et $g''$ s'annule sans changer de signe sur $\R$.
	
La fonction $g$ est donc convexe sur $\R$. \hfill \fbox{\textbf{Réponse b}}
	\item D'après la représentation graphique de $f'$, on peut affirmer que $f'(0)=1$. La pente de la tangente à $\mathcal{C}$ au point d'abscisse 0 est égale à 1. Toute droite ayant la même pente est donc parallèle à cette tangente. C'est le cas de la droite d'équation $y=x$. 
	
\hfill \fbox{\textbf{Réponse a}}
	\item On considère la suite $\left(u_n\right)$ définie pour tout entier naturel $n$ par $u_n = \dfrac{(-1)^n}{n+1}$.
	
$\forall n \in \N$, $ -1\leqslant (-1)^n \leqslant 1$ donc $-\dfrac{1}{n+1} \leqslant u_n \leqslant \dfrac{1}{n+1}$.
	
Or $\forall n \in \N$, $n+1\geqslant 1$ donc $\dfrac{1}{n+1} \leqslant 1$ et $-\dfrac{1}{n+1} \geqslant -1$.
	
Donc  $\forall n \in \N$, $-1\leqslant u_n \leqslant 1$. La suite $(u_n)$ est donc bornée.
	
\hfill \fbox{\textbf{Réponse c}}
	\item Soit $\left(v_n\right)$ une suite telle que $v_0 = k$ et $\forall n \in \N$, $v_n\times v_{n+1}<0$. Cette inégalité nous permet d'affirmer que $\forall n \in \N$,  deux termes consécutifs $v_n$ et $v_{n+1}$ sont de signes opposés. Donc les termes  $v_{n+1}$ et  $v_{n+2}$ le sont aussi. Donc on peut en déduire que $\forall n \in \N$, $v_n$ et $v_{n+2}$  sont de même signe.
	
Donc $v_0$, $v_2$, ..., $v_{2k}$, \: avec $k\in \N$ (tous les termes de rangs pairs), sont de même signe, donc du signe de $k$.
	
\hfill \fbox{\textbf{Réponse c}}
	\item On considère la suite $\left(w_n\right)$ définie pour tout entier naturel $n$ par $w_{n+1} = 2w_n - 4\quad \text{et}\quad  w_2 = 8$.
	
Donc $\forall n \in \N$,  $w_n=\dfrac{w_{n+1}+4}{2}$.
	
$w_1=\dfrac{w_{2}+4}{2}=\dfrac{8+4}{2}=6$ et $w_0=\dfrac{w_{1}+4}{2}=\dfrac{6+4}{2}=5$ \hfill \fbox{\textbf{Réponse b}}
	\item Il est facile de démontrer par récurrence que $\forall n \in \N$, $a_n > 0$ car $\dfrac{\e^n}{\e^n+1}> 0 $ et $a_0 > 0$.
	
$\forall n \in \N$, $\e^n <  \e^n+1$ donc $\dfrac{\e^n}{\e^n+1}< 1$ donc $\dfrac{\e^n}{\e^n+1} a_n< a_n$ donc $a_{n+1} < a_n$. La suite $(a_n)$ est donc strictement décroissante. \hfill \fbox{\textbf{Réponse b}}
	
	\item D'après l'énoncé, nous savons que le nombre de cellules double à chaque intervalle de temps écoulé. Cherchons le premier entier $n$ tel que $2^n\geqslant \np{4000}$. La fonction $x\mapsto \ln(x)$ est strictement croissante sur $\R_+^*$ donc 
	
$2^n\geqslant \np{4000}\iff \ln\left(2^n\right) \geqslant \ln(\np{4000})\iff n \times \ln(2) \geqslant \ln(\np{4000}) \iff n \geqslant \dfrac{\ln(\np{4000}) }{\ln(2)}$.
	
À la calculatrice : $\dfrac{\ln(\np{4000}) }{\ln(2)}\approx 11,97$ donc $n\geqslant 12.$
	
Il s'est donc écoulé 12 intervalles de temps pour que le nombre de cellules atteigne $\np{4000}$ en 4 heures. Chaque intervalle de temps est donc : $\dfrac{4\times 60}{12}= 20$ (min.).\hfill \fbox{\textbf{Réponse c}}
	
\end{enumerate}

\bigskip

\end{document}
